\documentclass[11pt,a4paper]{article}
\usepackage[T1]{fontenc}
\usepackage[utf8]{inputenc}
\usepackage{tgtermes}
\usepackage[margin=2.2cm,top=2.2cm,bottom=2.4cm]{geometry}
\usepackage{graphicx,booktabs,xcolor,caption,amsmath,float,array}
\usepackage{microtype}
\usepackage[colorlinks=true]{hyperref}

\definecolor{ecink}{HTML}{1A2430}
\definecolor{ecgray}{HTML}{5A6572}
\definecolor{citeblue}{HTML}{2A5DB0}
\hypersetup{linkcolor=citeblue,urlcolor=citeblue}
\captionsetup{font={small},labelsep=period,
              justification=justified,singlelinecheck=false}
\setlength{\parskip}{5pt}\setlength{\parindent}{0pt}
\color{ecink}
\newcommand{\fignote}[1]{\par\vspace{4pt}%
  \begin{minipage}{\linewidth}\footnotesize\color{ecgray}#1\end{minipage}}

\begin{document}

\begin{center}
{\fontsize{15}{20}\selectfont\bfseries
Teacher Attrition in the Seven CPF Countries\par}
\vspace{4pt}
{\fontsize{11.5}{15}\selectfont
How Often Do Teachers Leave the Profession, and How Do We Know?\par}
\vspace{8pt}
{\small José Elías Durán Roa \quad\textcolor{ecgray}{University of
Alicante and OECD} \quad\textcolor{ecgray}{July 2026}}
\end{center}
\vspace{4pt}

The Comparative Panel File (CPF) harmonizes the seven longest-running
national household panels: the PSID (United States), BHPS/Understanding
Society (United Kingdom), SOEP (Germany), HILDA (Australia), KLIPS
(South Korea), RLMS-HSE (Russia), and SHP (Switzerland). For each of
the seven countries, this note assembles the annual rate at which
school teachers leave the teaching profession, contrasting at least
three independent within-country sources per country---administrative
censuses and registers, household and follow-up surveys, and academic
studies---about twenty-five sources in all. Table~\ref{tab:head}
summarizes the country-level answer; Table~\ref{tab:detail} lists the
underlying sources.

Three regularities organize the evidence. First, within every country
the measured rate falls along the same definitional ladder: gross
separations from individual schools exceed exits from the school
sector, which exceed exits from the profession, which exceed
definitive exits net of returns, which exceed voluntary quits. Most
of the apparent disagreement between sources disappears once
definitions are aligned; in Switzerland, for example, the 8.4 percent
of teachers who leave the school registers each year shrinks to 1--2
percent once retirements and the 61 percent of under-55 leavers who
return within four years are netted out. Second, attrition scales
inversely with civil-service protection: voluntary quits are nearly
nil where teachers are tenured civil servants (Korea 0.2--0.3 percent,
German Beamte 0.2 percent) and the all-cause rate rises from about 2
percent in Korea to 5 percent in Australia, 4--6 percent in Russia,
7.5--8 percent in the United States, and roughly 9 percent leaving the
state-funded sector in England. The United States sits at the high end
of the distribution but is no outlier. Third, the household panels of
the CPF themselves are largely unable to measure teacher attrition:
no published teacher-exit rate exists for the PSID, SOEP, KLIPS,
RLMS-HSE, or SHP, whose per-wave teacher samples are too small, and
where household-panel estimates do exist they run roughly double the
administrative benchmark (12.3 percent in Understanding Society
against 9 percent in the English workforce census; 11.6 percent in
HILDA against 5 percent in Australian tax records), a gap that
Australian work traces to occupational miscoding between waves. This
is precisely the measurement problem the March CPS design
avoids---one instrument, annual, public, validated person by person
against panel classifications---and it is why the retrospective
instrument, not the household panels, carries the burden of
measurement in this literature. Everywhere, finally, stated intentions
vastly overstate behavior: 36--39 percent of Australian teachers
intend to leave and about 5 percent do; 87 percent of Korean teachers
have considered leaving and 0.2--0.3 percent resign.

\begin{table}[H]
\centering\small
\caption{Annual Teacher Profession-Leaving Rates in the CPF Countries}
\label{tab:head}
\begin{tabular}{lccc}
\toprule
Country (CPF panel) & Best estimate & Years & Range across sources \\
\midrule
United States (PSID)        & 7.5--8            & 2000s--2024      & 5--13.5 \\
United Kingdom (UKHLS)      & $\sim$9$^{a}$     & 2022/23--2024/25 & 5--12.3 \\
Germany (SOEP)              & 5--5.5$^{b}$      & 2010/11--2024/25 & 0.2--11.8 \\
Australia (HILDA)           & $\sim$5$^{c}$     & 2009--2021       & 1.3--11.6 \\
South Korea (KLIPS)         & $\sim$2$^{d}$     & 2019--2024       & 0.1--2.1 \\
Russia (RLMS-HSE)           & 4--6$^{e}$        & 2016--2024       & $<$1--14.4 \\
Switzerland (SHP)           & 7--8$^{f}$        & 2010--2021       & 1--8.4 \\
\bottomrule
\end{tabular}
\fignote{\textit{Notes:} Percent of the teacher stock leaving the
teaching profession per year; ``range'' spans all credible
within-country sources under their own definitions.
$^{a}$Share leaving the state-funded sector (England); genuine
profession exit is an estimated 5--8 percent.
$^{b}$Permanent exits from the school service including regular
retirement; voluntary quits are 0.2--0.4 percent.
$^{c}$All-cause; 4.5 percent among working-age teachers (21--64).
$^{d}$Departures before statutory retirement age, dominated by
honorary early retirement; voluntary resignations are 0.2--0.3
percent.
$^{e}$Early-career cohort exit is firmly documented at about 6
percent; official gross school separations (14.4 percent in 2023)
conflate profession exit with school-to-school moves.
$^{f}$Gross annual exits from the school registers, of which about 3
points are retirements; definitive profession-leaving of under-55s is
1--2 percent after re-entries.
\textit{Source:} Author's compilation, July 2026; sources in
Table~\ref{tab:detail}.}
\end{table}

\begin{table}[H]
\centering\footnotesize
\setlength{\tabcolsep}{4pt}
\caption{Underlying Sources, by Country}
\label{tab:detail}
\begin{tabular}{p{2.2cm}p{4.4cm}p{4.2cm}p{1.9cm}p{2.5cm}}
\toprule
Country & Source / dataset & Definition & Years & Annual rate (\%) \\
\midrule
United States & NCES Teacher Follow-up Survey (TFS/NTPS) & left K--12 teaching, public & 1988--2022 & 5.6 $\rightarrow$ 7.4--8.4; 8.0 (2021--22) \\
 & March CPS: Harris--Adams; Aldeman--Yi & occupation exit, all sectors & 1992--2001; 2015--24 & 7.7; 7.6 \\
 & State admin.: Washington S-275; Texas PEIMS & leaving state teaching workforce & 1985--2025 & 5--8 (WA); 9--13.5 (TX) \\
 & Learning Policy Institute (SASS/TFS) & leavers & 2011--13 & $\sim$8 \\
\addlinespace
United Kingdom & DfE School Workforce Census (England) & leaving state-funded sector & 2010/11--24/25 & 8.5--10.6 \\
 & NFER (SWC; excl.\ retirements) & working-age sector exit & 2022/23 & 8.8 \\
 & NFER on Understanding Society & profession exit, panel & 2009--17 & 12.3 (upper bound) \\
 & GTC Scotland; EWC Wales registers & de-registration & 2024--25 & $\sim$5; $\sim$4.5--5 \\
\addlinespace
Germany & Destatis Fachserie 11 / FiBS & permanent exits from school service & 2010/11--24/25 & $\sim$5.4 (gross 9.5--11.8) \\
 & Destatis pension statistics & new teacher retirements (Beamte) & 2003--23 & $\sim$2--3 \\
 & L\"ander data (NRW, BW) & voluntary resignation & 2013--24 & 0.2 (Beamte); 2--2.4 (angestellte) \\
\addlinespace
Australia & e61 Institute, PLIDA/BLADE tax data & not teaching in any of next 3 years & 2009--21 & 5.9 $\rightarrow$ 5.1; 4.5 working-age \\
 & AITSL ATWD registration data & registration discontinuation & 2019--23 & $\sim$1.3 (early-career) \\
 & HILDA (Cuervo--Vera-Toscano) & occupation exit, panel & 2001--22 & 11.6 (upper bound) \\
 & NSW CESE; Queensland College of Teachers & employer separations; register removal & 2006--19 & 3.7--5.3; $\sim$3.6 \\
\addlinespace
South Korea & MOE/KEDI administrative data & pre-retirement departures; resignations & 2019--24 & 1.5--2.1; 0.2--0.3 \\
 & Jung et al.\ (2021), provincial panel & profession exit, public elementary & 2009--18 & $<$2 \\
 & OECD TALIS 2024 (KEDI) & 5-year leaving intention & 2024 & $\le$1.2 implied \\
\addlinespace
Russia & Minprosveshcheniya form OO-1 & school separations (incl.\ moves) & 2017--23 & 10.7 $\rightarrow$ 14.4 (upper bound) \\
 & HSE graduate-employment panels & school-system exit, early career & 2016--24 & 5.9--6.9 \\
 & Tochno.st (OO-1/Rosstat) & net stock change; education-sector exit & 2016--25 & $-$0.9 net; $\sim$2.4 \\
\addlinespace
Switzerland & BFS LABB, Mobilit\"at (2014) & exit from school service & 2010--11 & 8.4; 16 first-year \\
 & BFS LABB, Verbleib (2022) & under-55 school-employment exit & 2015/16--20/21 & $\sim$2 ($\sim$1 definitive) \\
 & Canton Z\"urich; Herzog--Sandmeier; SKBF & system exit $\le$54; profession exits; retirements & 2016--24 & $\sim$7; 7--8; 3 \\
\bottomrule
\end{tabular}
\fignote{\textit{Notes:} One row per source family; rates as reported
by each source under its own definition, in percent of the relevant
teacher stock per year. Arrows denote change over the period.
``Upper bound'' marks household-panel estimates inflated by
occupational coding churn between waves. Intention measures are
excluded except the TALIS row, shown as an implied upper bound. For
the PSID, SOEP, KLIPS, RLMS-HSE, and SHP no published teacher-exit
rate exists; those countries rely on administrative sources. Full
citations and URLs in the replication file
\texttt{outputs/cpf\_teacher\_attrition.md}.
\textit{Source:} Author's compilation from the sources listed, July
2026.}
\end{table}

\end{document}
